\documentclass[11pt]{article}
\usepackage{amsmath,amssymb,amsthm}
\usepackage{mathrsfs}
\usepackage[margin=1in]{geometry}

\newtheorem{definition}{Definition}[section]
\newtheorem{conjecture}{Conjecture}[section]
\newtheorem{theorem}{Theorem}[section]
\newtheorem{lemma}[theorem]{Lemma}
\newtheorem{proposition}[theorem]{Proposition}
\newtheorem{corollary}[theorem]{Corollary}
\newtheorem{remark}{Remark}[section]

\newcommand{\R}{\mathbb{R}}
\newcommand{\E}{\mathbb{E}}
\newcommand{\Prob}{\mathbb{P}}
\newcommand{\calH}{\mathcal{H}}
\newcommand{\calM}{\mathcal{M}}
\newcommand{\calO}{\mathcal{O}}
\newcommand{\calS}{\mathcal{S}}
\newcommand{\supp}{\mathrm{supp}}
\newcommand{\dH}{d_{\mathrm{H}}}
\newcommand{\Wp}{W_p}

\title{Solution-Manifold Universality: \\Formal Conjectures and Proof Attempts}
\author{}
\date{\today}

\begin{document}
\maketitle

\begin{abstract}
We formalize four conjectures arising from the hypothesis that universality in stochastic dynamics is fundamentally about attractor geometry rather than parameter equality. We provide rigorous definitions, state the conjectures precisely, and give proof attempts—identifying what can be proved, what requires additional assumptions, and what remains open.
\end{abstract}

\tableofcontents

%==============================================================================
\section{Preliminaries and Definitions}
%==============================================================================

\subsection{Stochastic Dynamics and Solution Spaces}

\begin{definition}[Stochastic Growth Process]
A \textbf{stochastic growth process} on domain $\Lambda_L = [0,L]^d$ with periodic boundary conditions is a probability measure $\Prob$ on the path space
\[
\calH_{L,T} = C([0,T]; H^s(\Lambda_L))
\]
for appropriate Sobolev index $s$, satisfying a stochastic PDE of the form
\[
\partial_t h = \mathcal{L}[h] + \eta
\]
where $\mathcal{L}$ is a (possibly nonlinear) differential operator and $\eta$ is space-time white noise.
\end{definition}

\begin{definition}[Solution Manifold]
For a stochastic process with stationary measure $\mu$ on $H^s(\Lambda_L)$, the \textbf{solution manifold} at scale $L$ is
\[
\calM_L := \supp(\mu) \subset H^s(\Lambda_L)
\]
When $\mu$ concentrates on a lower-dimensional set (in some appropriate sense), we say the dynamics has a \textbf{finite-dimensional attractor}.
\end{definition}

\subsection{Observable Maps and Coarse-Graining}

\begin{definition}[Observable Map]
An \textbf{observable map} is a measurable function $\Phi: H^s(\Lambda_L) \to \R^k$ that extracts finite-dimensional statistics from height configurations. The \textbf{induced measure} is
\[
\mu^\Phi := \Phi_* \mu = \mu \circ \Phi^{-1}
\]
\end{definition}

\begin{definition}[Coarse-Graining Operator]
A \textbf{coarse-graining operator} at scale $\ell$ is a map $C_\ell: H^s(\Lambda_L) \to H^s(\Lambda_{L/\ell})$ that averages over regions of size $\ell$. For a height field $h$:
\[
(C_\ell h)(x) = \ell^{-\alpha} \int_{|y| < \ell/2} h(x + y) \, dy
\]
where $\alpha$ is chosen for stationarity of the limit.
\end{definition}

\begin{definition}[Local Differential Invariants]
The \textbf{local differential invariants} of order $n$ are the symmetric functions of derivatives up to order $n$:
\[
\mathcal{I}^{(n)} = \{\nabla h, \nabla^2 h, \ldots, \nabla^n h\}
\]
The \textbf{local invariant distribution} is the joint law of $\mathcal{I}^{(n)}(x)$ at a typical point $x$ under the stationary measure.
\end{definition}

\subsection{Distances Between Measures}

\begin{definition}[Wasserstein Distance]
For probability measures $\mu, \nu$ on $\R^k$ with finite $p$-th moments, the \textbf{Wasserstein-$p$ distance} is
\[
\Wp(\mu, \nu) = \left( \inf_{\gamma \in \Gamma(\mu,\nu)} \int \|x - y\|^p \, d\gamma(x,y) \right)^{1/p}
\]
where $\Gamma(\mu,\nu)$ is the set of couplings with marginals $\mu, \nu$.
\end{definition}

\begin{definition}[Hausdorff Distance on Supports]
For compact sets $A, B \subset \R^k$, the \textbf{Hausdorff distance} is
\[
\dH(A, B) = \max\left\{ \sup_{a \in A} \inf_{b \in B} \|a-b\|, \sup_{b \in B} \inf_{a \in A} \|a-b\| \right\}
\]
\end{definition}

%==============================================================================
\section{The Four Core Conjectures}
%==============================================================================

\subsection{Conjecture A: Attractor Equivalence Implies Exponent Equality}

\begin{conjecture}[Attractor $\Rightarrow$ Exponents]
\label{conj:attractor-exponents}
Let $\mathcal{L}_1, \mathcal{L}_2$ be two growth dynamics with stationary measures $\mu_1, \mu_2$ on $H^s(\Lambda_L)$. Let $\Phi: H^s \to \R^k$ be the standard observable map (interface width, height-height correlation, etc.). If there exists a measure-preserving homeomorphism
\[
\psi: \supp(\mu_1^\Phi) \to \supp(\mu_2^\Phi)
\]
then the scaling exponents coincide: $\alpha_1 = \alpha_2$, $\beta_1 = \beta_2$, $z_1 = z_2$.
\end{conjecture}

\begin{remark}
This conjecture asserts that attractor geometry determines exponents. The converse is known to be false (BD has same $\alpha$ as KPZ but different attractor).
\end{remark}

\textbf{Proof Attempt:}

\begin{proof}[Partial Proof under Additional Assumptions]
Assume:
\begin{enumerate}
\item[(A1)] The observable $\Phi$ includes the interface width: $w(h) = \sqrt{\frac{1}{L}\int (h - \bar{h})^2}$
\item[(A2)] The homeomorphism $\psi$ is bi-Lipschitz with constant $K$: $K^{-1}\|x-y\| \leq \|\psi(x) - \psi(y)\| \leq K\|x-y\|$
\item[(A3)] The scaling $w \sim L^\alpha$ holds in the sense that $\mu^\Phi$ concentrates near the curve $\{w = c \cdot L^\alpha\}$
\end{enumerate}

Under (A1)-(A3): The width $w$ defines a coordinate on $\supp(\mu^\Phi)$. The measure $\mu_1^\Phi$ concentrates near $w_1 = c_1 L^{\alpha_1}$. Under $\psi$, this maps to a region of $\supp(\mu_2^\Phi)$.

By (A2), the image has diameter at most $K \cdot \delta_1$ where $\delta_1$ is the concentration diameter of $\mu_1^\Phi$. 

If $\alpha_1 \neq \alpha_2$, then for large $L$:
\[
|w_1 - w_2| = |c_1 L^{\alpha_1} - c_2 L^{\alpha_2}| \to \infty
\]
But measure-preservation requires $\psi_*\mu_1^\Phi = \mu_2^\Phi$, so the supports must have comparable locations. Contradiction.

Therefore $\alpha_1 = \alpha_2$.

\textbf{Gap:} This argument requires (A2), which is not guaranteed by mere homeomorphism. A homeomorphism can stretch arbitrarily. The conjecture may need strengthening to ``bi-Lipschitz equivalence'' or weakening to ``scaling equivalence up to logarithmic corrections.''
\end{proof}

\begin{proposition}[Weaker Version]
If $\psi: \supp(\mu_1^\Phi) \to \supp(\mu_2^\Phi)$ is a \textbf{bi-Lipschitz} measure-preserving map with constant $K$ independent of $L$, then $\alpha_1 = \alpha_2$.
\end{proposition}

\begin{proof}
The concentration diameters satisfy $\delta_i(L) \sim L^{-\gamma_i}$ for some $\gamma_i > 0$. Under bi-Lipschitz map:
\[
K^{-1} \delta_1(L) \leq \delta_2(L) \leq K \delta_1(L)
\]
Taking logs: $-\gamma_1 \log L + O(1) \approx -\gamma_2 \log L + O(1)$, giving $\gamma_1 = \gamma_2$.

The exponent $\alpha$ is related to $\gamma$ through the Family-Vicsek scaling. For EW, $\gamma = 1/2$ and $\alpha = 1/2$. For KPZ, $\gamma \approx 1/6$ (heuristic) and $\alpha = 1/2$. 

\textbf{Issue:} $\gamma$ depends on the observable, not just $\alpha$. The connection $\gamma \leftrightarrow \alpha$ requires the specific form of $\Phi$. This is the ``relevant observable'' problem.
\end{proof}

%------------------------------------------------------------------------------
\subsection{Conjecture B: Local Constraints Determine Universality}

\begin{conjecture}[Local Carrier Hypothesis]
\label{conj:local-carrier}
For stochastic growth processes with local interactions, universality class membership is determined by the joint distribution of local differential invariants. Precisely: if $\mu_1^{\mathcal{I}^{(n)}} = \mu_2^{\mathcal{I}^{(n)}}$ for some finite $n$, then $\mathcal{L}_1$ and $\mathcal{L}_2$ are in the same universality class.
\end{conjecture}

\textbf{Proof Attempt:}

\begin{proof}[Proof for EW vs KPZ at $n=1$]
Consider the Edwards-Wilkinson and KPZ equations:
\begin{align}
\text{EW:} \quad \partial_t h &= \nu \nabla^2 h + \eta \\
\text{KPZ:} \quad \partial_t h &= \nu \nabla^2 h + \frac{\lambda}{2}(\nabla h)^2 + \eta
\end{align}

\textbf{Step 1: Gradient distribution under EW.}

EW is a linear SDE. The stationary measure is Gaussian with covariance $\langle h(x)h(y)\rangle \propto |x-y|^{2\alpha}$ where $\alpha = 1/2$ in 1+1D. The gradient $\nabla h$ is also Gaussian:
\[
\nabla h \sim \mathcal{N}(0, \sigma^2_{\text{EW}})
\]
where $\sigma^2_{\text{EW}} = D/\nu$ depends on noise strength $D$ and viscosity $\nu$.

\textbf{Step 2: Gradient distribution under KPZ.}

The Cole-Hopf transformation $Z = e^{\lambda h / 2\nu}$ maps KPZ to the stochastic heat equation:
\[
\partial_t Z = \nu \nabla^2 Z + \frac{\lambda}{2\nu} Z \eta
\]
The gradient $\nabla h = (2\nu/\lambda) \nabla \log Z$ has a \textbf{non-Gaussian} stationary distribution. 

Specifically, the one-point distribution of $\nabla h$ under KPZ at stationarity has nonzero skewness. For the KPZ fixed point (Matetski-Quastel-Remenik), the gradient distribution is related to the Airy process derivative, which is asymmetric.

\textbf{Step 3: Separation.}

Define:
\[
\gamma_1[\nabla h] := \E[(\nabla h)^3] / \E[(\nabla h)^2]^{3/2}
\]
For EW: $\gamma_1^{\text{EW}} = 0$ (Gaussian symmetry). For KPZ: $\gamma_1^{\text{KPZ}} \neq 0$ (nonlinear term breaks symmetry).

Therefore $\mu_{\text{EW}}^{\nabla h} \neq \mu_{\text{KPZ}}^{\nabla h}$, and this is detectable at $n=1$.

\textbf{Conclusion:} EW and KPZ are distinguished by their first-order local invariant distribution. \qed
\end{proof}

\begin{theorem}[Gradient Skewness Separation]
\label{thm:gradient-skewness}
Let $\mu_{\text{EW}}, \mu_{\text{KPZ}}$ be the stationary measures for EW and KPZ on $[0,L]$ with periodic boundaries. Let $\Phi(h) = \gamma_1[\nabla h]$ be the sample skewness of the gradient. Then:
\[
\lim_{L \to \infty} W_1(\mu_{\text{EW}}^\Phi, \mu_{\text{KPZ}}^\Phi) > 0
\]
\end{theorem}

\begin{proof}
By the argument above:
\begin{itemize}
\item $\mu_{\text{EW}}^\Phi$ concentrates at $\gamma_1 = 0$ (exact, from Gaussianity)
\item $\mu_{\text{KPZ}}^\Phi$ concentrates at $\gamma_1 = \gamma_* \neq 0$ (from KPZ fixed point asymptotics)
\end{itemize}

The concentration follows from spatial CLT: sample skewness from $L$ points has variance $O(1/L)$.

Therefore $W_1 \geq |\gamma_*| - O(L^{-1/2}) > 0$ for large $L$.
\end{proof}

\begin{remark}
This is a ``gradient version'' of Theorem 5 in main.tex (which uses height skewness). The gradient version may have better finite-size behavior because gradient fluctuations are more localized.
\end{remark}

\textbf{Open Question:} Is there a converse? If $\mu_1^{\mathcal{I}^{(n)}} = \mu_2^{\mathcal{I}^{(n)}}$ for all $n$, does this imply same universality class? 

This would require showing that local invariant distributions are \textbf{complete invariants}---they determine the full stationary measure. This is likely false in general (non-local correlations matter), but may hold for translation-invariant SPDEs.

%------------------------------------------------------------------------------
\subsection{Conjecture C: Solver-Independent Distance}

\begin{conjecture}[Intrinsic Physics Distance]
\label{conj:solver-independent}
There exists a distance function $d_{\text{phys}}$ on the space of stochastic growth laws such that:
\begin{enumerate}
\item $d_{\text{phys}}(\mathcal{L}_1, \mathcal{L}_2) = 0$ iff $\mathcal{L}_1, \mathcal{L}_2$ are in the same universality class
\item For any consistent numerical scheme $S$ with discretization $\epsilon$:
\[
d_{\text{phys}}^{S,\epsilon}(\mathcal{L}_1, \mathcal{L}_2) \to d_{\text{phys}}(\mathcal{L}_1, \mathcal{L}_2) \quad \text{as } \epsilon \to 0
\]
\item The convergence is uniform over compact families of growth laws
\end{enumerate}
\end{conjecture}

\textbf{Proposed Construction:}

\begin{definition}[Stationary Measure Distance]
For growth laws $\mathcal{L}_1, \mathcal{L}_2$ with stationary measures $\mu_1, \mu_2$ on gradient observables:
\[
d_{\text{phys}}(\mathcal{L}_1, \mathcal{L}_2) := \lim_{L \to \infty} W_1(\mu_1^{\nabla h}, \mu_2^{\nabla h})
\]
where $\mu_i^{\nabla h}$ is the marginal distribution of $\nabla h(0)$ under the stationary measure at system size $L$.
\end{definition}

\begin{proposition}[Property 1: Universality Metrization]
If $\mathcal{L}_1, \mathcal{L}_2$ are in the same universality class (flow to same RG fixed point), then $d_{\text{phys}}(\mathcal{L}_1, \mathcal{L}_2) = 0$.
\end{proposition}

\begin{proof}[Heuristic Argument]
Under RG flow, the gradient distribution flows toward the fixed point distribution. If $\mathcal{L}_1, \mathcal{L}_2$ share a fixed point $\mathcal{L}_*$, then:
\[
\mu_1^{\nabla h} \to \mu_*^{\nabla h}, \quad \mu_2^{\nabla h} \to \mu_*^{\nabla h}
\]
as $L \to \infty$ (the infrared limit). Therefore:
\[
W_1(\mu_1^{\nabla h}, \mu_2^{\nabla h}) \leq W_1(\mu_1^{\nabla h}, \mu_*^{\nabla h}) + W_1(\mu_*^{\nabla h}, \mu_2^{\nabla h}) \to 0
\]

\textbf{Gap:} This requires the RG flow to converge in Wasserstein distance, not just in distribution. The Cotler-Rezchikov theorem establishes this for Euclidean QFT, but extension to stochastic PDEs is not proven.
\end{proof}

\begin{proposition}[Property 2: Numerical Convergence]
For consistent finite-difference schemes, $d_{\text{phys}}^{S,\epsilon} \to d_{\text{phys}}$ as $\epsilon \to 0$.
\end{proposition}

\begin{proof}[Sketch]
A consistent scheme $S$ with discretization $\epsilon$ produces a Markov chain whose stationary measure $\mu^{S,\epsilon}$ approximates the continuum stationary measure $\mu$. 

Standard results in numerical SPDEs give weak convergence: $\mu^{S,\epsilon} \Rightarrow \mu$ as $\epsilon \to 0$.

Weak convergence plus uniform integrability implies $W_1(\mu^{S,\epsilon}, \mu) \to 0$.

Therefore:
\begin{align}
|d_{\text{phys}}^{S,\epsilon} - d_{\text{phys}}| &= |W_1(\mu_1^{S,\epsilon}, \mu_2^{S,\epsilon}) - W_1(\mu_1, \mu_2)| \\
&\leq W_1(\mu_1^{S,\epsilon}, \mu_1) + W_1(\mu_2^{S,\epsilon}, \mu_2) \to 0
\end{align}

\textbf{Gap:} Uniform integrability must be verified for specific SPDEs. For KPZ, tail bounds on the stationary measure are known (Tracy-Widom tails), making this tractable.
\end{proof}

%------------------------------------------------------------------------------
\subsection{Conjecture D: Coarse-Grained Manifold Convergence}

\begin{conjecture}[Manifold Coalescence]
\label{conj:coalescence}
Let $\mathcal{L}_1, \mathcal{L}_2$ be in the same universality class. Let $\calM_1^{(L)}, \calM_2^{(L)}$ be their solution manifolds at scale $L$, projected to gradient observable space. Then:
\[
\dH(\calM_1^{(L)}, \calM_2^{(L)}) \to 0 \quad \text{as } L \to \infty
\]
That is, the solution manifolds \textbf{merge} under coarse-graining.
\end{conjecture}

\textbf{Proof Attempt:}

\begin{proof}[Proof under Concentration Hypothesis]
Assume:
\begin{itemize}
\item[(H1)] Both manifolds concentrate: $\text{diam}(\calM_i^{(L)}) \leq C \cdot L^{-\gamma}$ for some $\gamma > 0$
\item[(H2)] The centers converge: $\text{center}(\calM_1^{(L)}) \to m_*$, $\text{center}(\calM_2^{(L)}) \to m_*$
\end{itemize}

Under (H1) and (H2):
\[
\dH(\calM_1^{(L)}, \calM_2^{(L)}) \leq \|\text{center}_1 - \text{center}_2\| + \text{diam}(\calM_1) + \text{diam}(\calM_2)
\]
\[
\leq o(1) + C L^{-\gamma} + C L^{-\gamma} \to 0
\]

\textbf{Verification of (H1):} For EW, concentration at rate $L^{-1/2}$ is proved (CLT). For KPZ, $L^{-1/6}$ is heuristic (main.tex Appendix).

\textbf{Verification of (H2):} If both flow to the same RG fixed point, the limiting gradient distribution is identical, so centers converge to the same point.
\end{proof}

\begin{theorem}[EW Manifold Coalescence]
Let $\mathcal{L}_1 = $ EW with parameters $(\nu_1, D_1)$ and $\mathcal{L}_2 = $ EW with parameters $(\nu_2, D_2)$. Then:
\[
\dH(\calM_1^{(L)}, \calM_2^{(L)}) = O(L^{-1/2})
\]
\end{theorem}

\begin{proof}
Both are EW, hence same universality class. The stationary gradient distribution for EW$(\nu, D)$ is $\mathcal{N}(0, D/\nu)$.

After standardization (rescaling to unit variance), both have gradient distribution $\mathcal{N}(0, 1)$.

The sample gradient variance from $L$ points has distribution concentrated around mean $1$ with standard deviation $O(L^{-1/2})$.

Therefore manifolds (in standardized coordinates) have diameter $O(L^{-1/2})$ and both center at $(1, 0, \ldots)$ in the (variance, skewness, ...) coordinate system.

Hence $\dH = O(L^{-1/2})$.
\end{proof}

%==============================================================================
\section{Synthesis: Toward a Theorem}
%==============================================================================

The four conjectures suggest a unified structure. We now state a combined result that captures the essence, with clearly marked assumptions.

\begin{theorem}[Solution-Manifold Universality --- Conditional]
\label{thm:main}
Let $\mathcal{L}_1, \mathcal{L}_2$ be stochastic growth processes in 1+1 dimensions satisfying:
\begin{enumerate}
\item[(A1)] \textbf{Ergodicity}: Both have unique ergodic stationary measures $\mu_1, \mu_2$
\item[(A2)] \textbf{Concentration}: The induced measures on gradient variance satisfy $\delta_i(L) \to 0$
\item[(A3)] \textbf{Local interaction}: The dynamics depend only on $h$ and its derivatives up to order $k$
\end{enumerate}

Then the following are equivalent:
\begin{enumerate}
\item[(i)] $\mathcal{L}_1, \mathcal{L}_2$ are in the same universality class (same RG fixed point)
\item[(ii)] $\lim_{L \to \infty} W_1(\mu_1^{\nabla h}, \mu_2^{\nabla h}) = 0$
\item[(iii)] $\lim_{L \to \infty} \dH(\calM_1^{(L)}, \calM_2^{(L)}) = 0$ in gradient observable space
\end{enumerate}
\end{theorem}

\begin{proof}[Proof Sketch]
\textbf{(i) $\Rightarrow$ (ii):} Same fixed point implies gradient distributions converge to the same limit (fixed point distribution). Wasserstein metrizes weak convergence for measures with bounded moments, giving $W_1 \to 0$.

\textbf{(ii) $\Rightarrow$ (iii):} $W_1 \to 0$ implies supports converge in Hausdorff distance (for concentrated measures).

\textbf{(iii) $\Rightarrow$ (i):} This is the hard direction. We argue by contrapositive.

Suppose $\mathcal{L}_1, \mathcal{L}_2$ have different fixed points $\mathcal{L}_1^*, \mathcal{L}_2^*$. Different fixed points have different gradient distributions (otherwise RG would identify them). By Conjecture~\ref{conj:local-carrier}, different gradient distributions are distinguished by some moment.

Let $\Phi_m$ be the discriminating moment (e.g., skewness if fixed points are EW vs KPZ). Then:
\[
|\E_{\mu_1^*}[\Phi_m] - \E_{\mu_2^*}[\Phi_m]| = \Delta > 0
\]

As $L \to \infty$, $\mu_i^{(L)} \to \mu_i^*$ in distribution, so:
\[
\E_{\mu_i^{(L)}}[\Phi_m] \to \E_{\mu_i^*}[\Phi_m]
\]

For large $L$:
\[
\dH(\calM_1^{(L)}, \calM_2^{(L)}) \geq |\E_{\mu_1}[\Phi_m] - \E_{\mu_2}[\Phi_m]| - (\delta_1 + \delta_2) \geq \Delta/2 > 0
\]

Contradiction with (iii). Therefore (iii) $\Rightarrow$ (i).
\end{proof}

\begin{remark}[Status of the Proof]
The argument for (i) $\Rightarrow$ (ii) $\Rightarrow$ (iii) is essentially complete given standard SPDE results.

The argument for (iii) $\Rightarrow$ (i) relies on:
\begin{itemize}
\item Different fixed points having different gradient distributions (Conjecture~\ref{conj:local-carrier})
\item RG flow converging in a suitable sense
\end{itemize}

A fully rigorous proof requires extending Cotler-Rezchikov's RG-Wasserstein theorem from Euclidean QFT to stochastic PDEs.
\end{remark}

%==============================================================================
\section{What Remains Open}
%==============================================================================

\begin{enumerate}
\item \textbf{Cotler-Rezchikov extension}: Prove that RG = Wasserstein gradient flow for stochastic PDEs, not just Euclidean QFT.

\item \textbf{Local completeness}: Characterize which universality classes are distinguished by gradient distributions alone, versus requiring higher-order invariants.

\item \textbf{Concentration rates}: Prove KPZ concentration at rate $L^{-1/6}$ (currently heuristic).

\item \textbf{Dimension bounds}: Prove that solution manifolds have finite intrinsic dimension, and characterize how this dimension relates to universality class.

\item \textbf{Beyond 1+1D}: Extend to higher dimensions where KPZ is less understood.
\end{enumerate}

%==============================================================================
\section{Empirical Predictions}
%==============================================================================

The conjectures make testable predictions:

\begin{enumerate}
\item \textbf{Gradient skewness test}: For any KPZ-class model (RSOS, ballistic deposition, PNG), the gradient skewness should converge to the same nonzero value.

\item \textbf{Intrinsic dimension}: Solution manifolds in gradient observable space should have $d_{\text{int}} = O(1)$, independent of system size.

\item \textbf{Solver convergence}: Different numerical schemes for the same SPDE should give converging estimates of $W_1(\text{EW}, \text{KPZ})$.

\item \textbf{Higher-order universality}: MBE ($\nabla^4$ equation) should be distinguished from EW/KPZ at $n=2$ (curvature statistics), not $n=1$.
\end{enumerate}

\end{document}
